%!TEX TS-program = xelatex
\documentclass[11pt]{article}

\usepackage[english]{babel}

\usepackage{amsmath,amssymb,amsfonts}
\usepackage[utf8]{inputenc}
\usepackage[T1]{fontenc}
\usepackage{stix}
\usepackage[scaled]{helvet}
\usepackage[scaled]{inconsolata}

\usepackage{lastpage}

\usepackage{setspace}

\usepackage{ccicons}

\usepackage[hang,flushmargin]{footmisc}

\usepackage{geometry}

\setlength{\parindent}{0pt}
\setlength{\parskip}{6pt plus 2pt minus 1pt}

\usepackage{fancyhdr}
\renewcommand{\headrulewidth}{0pt}\providecommand{\tightlist}{%
  \setlength{\itemsep}{0pt}\setlength{\parskip}{0pt}}

\makeatletter
\newcounter{tableno}
\newenvironment{tablenos:no-prefix-table-caption}{
  \caption@ifcompatibility{}{
    \let\oldthetable\thetable
    \let\oldtheHtable\theHtable
    \renewcommand{\thetable}{tableno:\thetableno}
    \renewcommand{\theHtable}{tableno:\thetableno}
    \stepcounter{tableno}
    \captionsetup{labelformat=empty}
  }
}{
  \caption@ifcompatibility{}{
    \captionsetup{labelformat=default}
    \let\thetable\oldthetable
    \let\theHtable\oldtheHtable
    \addtocounter{table}{-1}
  }
}
\makeatother

\usepackage{array}
\newcommand{\PreserveBackslash}[1]{\let\temp=\\#1\let\\=\temp}
\let\PBS=\PreserveBackslash

\usepackage[breaklinks=true]{hyperref}
\hypersetup{colorlinks,%
citecolor=blue,%
filecolor=blue,%
linkcolor=blue,%
urlcolor=blue}
\usepackage{url}

\usepackage{caption}
\setcounter{secnumdepth}{0}
\usepackage{cleveref}

\usepackage{graphicx}
\makeatletter
\def\maxwidth{\ifdim\Gin@nat@width>\linewidth\linewidth
\else\Gin@nat@width\fi}
\makeatother
\let\Oldincludegraphics\includegraphics
\renewcommand{\includegraphics}[1]{\Oldincludegraphics[width=\maxwidth]{#1}}

\usepackage{longtable}
\usepackage{booktabs}

\usepackage{color}
\usepackage{fancyvrb}
\newcommand{\VerbBar}{|}
\newcommand{\VERB}{\Verb[commandchars=\\\{\}]}
\DefineVerbatimEnvironment{Highlighting}{Verbatim}{commandchars=\\\{\}}
% Add ',fontsize=\small' for more characters per line
\usepackage{framed}
\definecolor{shadecolor}{RGB}{248,248,248}
\newenvironment{Shaded}{\begin{snugshade}}{\end{snugshade}}
\newcommand{\KeywordTok}[1]{\textcolor[rgb]{0.13,0.29,0.53}{\textbf{#1}}}
\newcommand{\DataTypeTok}[1]{\textcolor[rgb]{0.13,0.29,0.53}{#1}}
\newcommand{\DecValTok}[1]{\textcolor[rgb]{0.00,0.00,0.81}{#1}}
\newcommand{\BaseNTok}[1]{\textcolor[rgb]{0.00,0.00,0.81}{#1}}
\newcommand{\FloatTok}[1]{\textcolor[rgb]{0.00,0.00,0.81}{#1}}
\newcommand{\ConstantTok}[1]{\textcolor[rgb]{0.00,0.00,0.00}{#1}}
\newcommand{\CharTok}[1]{\textcolor[rgb]{0.31,0.60,0.02}{#1}}
\newcommand{\SpecialCharTok}[1]{\textcolor[rgb]{0.00,0.00,0.00}{#1}}
\newcommand{\StringTok}[1]{\textcolor[rgb]{0.31,0.60,0.02}{#1}}
\newcommand{\VerbatimStringTok}[1]{\textcolor[rgb]{0.31,0.60,0.02}{#1}}
\newcommand{\SpecialStringTok}[1]{\textcolor[rgb]{0.31,0.60,0.02}{#1}}
\newcommand{\ImportTok}[1]{#1}
\newcommand{\CommentTok}[1]{\textcolor[rgb]{0.56,0.35,0.01}{\textit{#1}}}
\newcommand{\DocumentationTok}[1]{\textcolor[rgb]{0.56,0.35,0.01}{\textbf{\textit{#1}}}}
\newcommand{\AnnotationTok}[1]{\textcolor[rgb]{0.56,0.35,0.01}{\textbf{\textit{#1}}}}
\newcommand{\CommentVarTok}[1]{\textcolor[rgb]{0.56,0.35,0.01}{\textbf{\textit{#1}}}}
\newcommand{\OtherTok}[1]{\textcolor[rgb]{0.56,0.35,0.01}{#1}}
\newcommand{\FunctionTok}[1]{\textcolor[rgb]{0.00,0.00,0.00}{#1}}
\newcommand{\VariableTok}[1]{\textcolor[rgb]{0.00,0.00,0.00}{#1}}
\newcommand{\ControlFlowTok}[1]{\textcolor[rgb]{0.13,0.29,0.53}{\textbf{#1}}}
\newcommand{\OperatorTok}[1]{\textcolor[rgb]{0.81,0.36,0.00}{\textbf{#1}}}
\newcommand{\BuiltInTok}[1]{#1}
\newcommand{\ExtensionTok}[1]{#1}
\newcommand{\PreprocessorTok}[1]{\textcolor[rgb]{0.56,0.35,0.01}{\textit{#1}}}
\newcommand{\AttributeTok}[1]{\textcolor[rgb]{0.77,0.63,0.00}{#1}}
\newcommand{\RegionMarkerTok}[1]{#1}
\newcommand{\InformationTok}[1]{\textcolor[rgb]{0.56,0.35,0.01}{\textbf{\textit{#1}}}}
\newcommand{\WarningTok}[1]{\textcolor[rgb]{0.56,0.35,0.01}{\textbf{\textit{#1}}}}
\newcommand{\AlertTok}[1]{\textcolor[rgb]{0.94,0.16,0.16}{#1}}
\newcommand{\ErrorTok}[1]{\textcolor[rgb]{0.64,0.00,0.00}{\textbf{#1}}}
\newcommand{\NormalTok}[1]{#1}

\newlength{\cslhangindent}
\setlength{\cslhangindent}{1.5em}
\newlength{\csllabelwidth}
\setlength{\csllabelwidth}{3em}
\newenvironment{CSLReferences}[3] % #1 hanging-ident, #2 entry spacing
 {% don't indent paragraphs
  \setlength{\parindent}{0pt}
  % turn on hanging indent if param 1 is 1
  \ifodd #1 \everypar{\setlength{\hangindent}{\cslhangindent}}\ignorespaces\fi
  % set entry spacing
  \ifnum #2 > 0
  \setlength{\parskip}{#2\baselineskip}
  \fi
 }%
 {}
\usepackage{calc} % for \widthof, \maxof
\newcommand{\CSLBlock}[1]{#1\hfill\break}
\newcommand{\CSLLeftMargin}[1]{\parbox[t]{\maxof{\widthof{#1}}{\csllabelwidth}}{#1}}
\newcommand{\CSLRightInline}[1]{\parbox[t]{\linewidth}{#1}}
\newcommand{\CSLIndent}[1]{\hspace{\cslhangindent}#1}\geometry{verbose,letterpaper,tmargin=2.5cm,bmargin=2.5cm,lmargin=2.5cm,rmargin=4.5cm}

\usepackage{lineno}
\usepackage[nolists,noheads]{endfloat}

\pagestyle{plain}

\doublespacing

\fancypagestyle{normal}
{
  \fancyhf{}
  \fancyfoot[R]{\footnotesize\sffamily\thepage\ of \pageref*{LastPage}}
}
\begin{document}
\thispagestyle{empty}
{\Large\bfseries\sffamily First Estimates of Food-Web Structure Derived
from Species Richness}
\vskip 5em

%
\href{https://orcid.org/0000-0001-9051-0597}{Francis\,Banville}%
%
\,\textsuperscript{1,2,3}\quad %
\href{https://orcid.org/0000-0002-4498-7076}{Dominique\,Gravel}%
%
\,\textsuperscript{2,3}\quad %
\href{https://orcid.org/0000-0002-0735-5184}{Timothée\,Poisot}%
%
\,\textsuperscript{1,3}

\textsuperscript{1}\,Université de
Montréal\quad \textsuperscript{2}\,Université de
Sherbrooke\quad \textsuperscript{3}\,Quebec Centre for Biodiversity
Science


\textbf{Correspondance to:}\\
Francis Banville --- \texttt{francis.banville@umontreal.ca}\\

\vfill
This work is released by its authors under a CC-BY 4.0 license\hfill\ccby\\
Last revision: \emph{\today}

\clearpage
\thispagestyle{empty}

\vfill
\vfill

\clearpage
\linenumbers
\pagestyle{normal}

\hypertarget{introduction}{%
\section{Introduction}\label{introduction}}

\hypertarget{methods}{%
\section{Methods}\label{methods}}

\hypertarget{a-maximum-entropy-model-for-predicting-food-web-structure}{%
\subsection{A maximum entropy model for predicting food-web
structure}\label{a-maximum-entropy-model-for-predicting-food-web-structure}}

Our mathematical model predicts the structure of species-based food webs
from their number of species. Our model involves three main steps: (i)
the prediction of the number of links from the number of species using
the flexible links model; (ii) the derivation of the joint degree
distribution using a maximum entropy approach; and (iii) the prediction
of the adjacency matrix using simulations and maximum entropy. The next
subsections describe each of these steps in detail.

\hypertarget{predicting-the-number-of-links}{%
\subsubsection{Predicting the number of
links}\label{predicting-the-number-of-links}}

In order to derive the joint degree distribution of maximum entropy, we
need two network-level measures: the number of species \(S\) and the
number of interactions \(L\). While the number of species is a
well-described measure of biodiversity for many taxa and locations, the
number of interactions is currently difficult to estimate empirically
without sampling all pairwise interactions in a biological community. We
thus used a predictive statistical model to simulate the number of
interactions from the number of species in food webs.

To do so, we used the flexible links model of MacDonald, Banville, and
Poisot (2020). The flexible links model incorporates meaningful
ecological constraints into the prediction of \(L\), namely the minimum
\(S-1\) and maximum \(S^2\) numbers of interactions in food webs, and
estimates the proportion of the \(S^2 - (S - 1)\) \emph{flexible links}
that are realized. More precisely, this model states that the number of
\emph{realized} flexible links \(L_{FL}\) in a food web represents the
number of realized interactions above the minimum
(i.e.~\(L = L_{FL} + S - 1\)) and is obtained from a beta-binomial
distribution with \(S^2 - (S - 1)\) trials and parameters
\(\alpha = \mu e^\phi\) and \(\beta = (1 - \mu) e^\phi\):

\begin{equation}\protect\hypertarget{eq:BB}{}{L_{FL} \sim \mathrm{BB}(S^2 - (S - 1), \mu e^\phi, (1 - \mu) e^\phi),}\label{eq:BB}\end{equation}

where \(\mu\) is the average probability across food webs that a
flexible link is realized, and \(\phi\) the concentration parameter
around \(\mu\).

We fitted the flexible links model on all food webs archived on Mangal,
an ecological interactions database (Poisot et al. 2016). Ecological
networks archived on Mangal are multilayer networks, i.e.~networks that
describe different types of interactions. We considered as food webs all
networks mainly composed of trophic interactions (predation and
herbivory types). We estimated the parameters of eq.~\ref{eq:BB} using a
Hamiltonian Monte Carlo sampler with static trajectory (1 chain and 3000
iterations):

\begin{equation}\protect\hypertarget{eq:BBpost}{}{
[\mu, \phi| \textbf{L}, \textbf{S}] \propto \prod_{i = 1}^{m} \mathrm{BB}(L_i - (S_i - 1) | S_i^2 - (S_i - 1)), \mu e^{\phi}, (1 - \mu) e^\phi) \times \mathrm{B}(\mu| 3 , 7 ) \times \mathcal{N}(\phi | 3, 0.5),
}\label{eq:BBpost}\end{equation}

where \(m\) is the number of food webs archived on Mangal (\(m = 235\))
and \(\textbf{L}\) and \(\textbf{S}\) are respectively the vectors of
their numbers of interactions and numbers of species. Our
weakly-informative prior distributions were chosen following MacDonald,
Banville, and Poisot (2020), i.e.~a beta distribution for \(\mu\) and a
normal distribution for \(\phi\). The Monte Carlo sampling of the
posterior distribution was conducted using the Julia library
\texttt{Turing} v0.15.12 (Ge, Xu, and Ghahramani 2018).

The flexible links model is a generative model, i.e.~it can generate
plausible values of the predicted variable. We thus simulated 1000
values of \(L\) for each value of \(S\) ranging between 5 and 1000
species using the joint posterior distribution of our model parameters.

\hypertarget{deriving-the-joint-degree-distribution}{%
\subsubsection{Deriving the joint degree
distribution}\label{deriving-the-joint-degree-distribution}}

The number of species and the number of interactions are state variables
whose ratio was directly used to obtain the least-biased degree
distribution \(p(k)\). This discrete probability distribution represents
the probability that a species have \(k\) interactions in its food web,
\(k\) ranging between 1 and \(S\). The least-biased of these
distributions, considering our limited knowledge of the system, can be
derived using the principle of maximum entropy, i.e.~by maximizing
Shannon's entropy

\begin{equation}\protect\hypertarget{eq:entropy}{}{H = -\sum_{k} p(k) \log p(k)}\label{eq:entropy}\end{equation}

while respecting a set of constraints on the degree distribution. We
used two such constraints:

\begin{equation}\protect\hypertarget{eq:g1}{}{g_1 = \sum_{k=1}^{S} p(k) = 1}\label{eq:g1}\end{equation}

and

\begin{equation}\protect\hypertarget{eq:g2}{}{g_2 = \langle k \rangle = \sum_{k=1}^{S} k p(k) = \frac{2L}{S}}\label{eq:g2}\end{equation}

The first constraint \(g_1\) is a normalizing constraint that ensures
that all probabilities sum to 1. In addition, the second constraint
\(g_2\) fixes the average of the degree distribution to the mean degree
\(\langle k \rangle\). The mean degree is twice the value of the linkage
density \(L/S\) since each interaction must be counted twice when
summing all species' degrees. We used each simulated value of \(L\) to
compute the mean degree constraint. As a result, we derived 1000 maximum
entropy degree distributions for each value of \(S\).

Finding probability distributions of maximum entropy is typically done
using the method of Lagrange multipliers:

\begin{equation}\protect\hypertarget{eq:lagrange1}{}{\frac{\partial H}{\partial p(k)} = \lambda_1 \frac{\partial g_1}{\partial p(k)} + \lambda_2 \frac{\partial g_2}{\partial p(k)},}\label{eq:lagrange1}\end{equation}

where \(\lambda_1\) and \(\lambda_2\) are the Lagrange multipliers.
Probability distributions of maximum entropy are derived by finding
these values. Evaluating the partial derivatives with respect to
\(p(k)\), we obtained:

\begin{equation}\protect\hypertarget{eq:lagrange2}{}{-\log p(k) - 1 = \lambda_1 + \lambda_2 k}\label{eq:lagrange2}\end{equation}

Then, solving eq.~\ref{eq:lagrange2} for \(p(k)\), we obtained

\begin{equation}\protect\hypertarget{eq:lagrange3}{}{p(k) = \frac{e^{-\lambda_2k}}{z},}\label{eq:lagrange3}\end{equation}

where \(z = e^{1+\lambda_1}\).

After substituting \(p(k)\) in eq.~\ref{eq:g1} and eq.~\ref{eq:g2}, we
got a nonlinear system of two equations and two unknowns:

\begin{equation}\protect\hypertarget{eq:lagrange4}{}{\frac{1}{z}\sum_{k=1}^{S}e^{-\lambda_2k} = 1}\label{eq:lagrange4}\end{equation}
\begin{equation}\protect\hypertarget{eq:lagrange5}{}{\frac{1}{z}\sum_{k=1}^{S}ke^{-\lambda_2k} = \frac{2L}{S},}\label{eq:lagrange5}\end{equation}

which implies that \(z = \sum_{k=1}^{S}e^{-\lambda_2k}\).

We solved eq.~\ref{eq:lagrange5} numerically using the Julia library
\texttt{NLsolve} v4.5.1 (TK). However, in figure TK, we show how an
analytical solution in the limit \(S \rightarrow \infty\) yields very
similar results. For large food webs (TK), the degree distribution of
maximum entropy approaches

\begin{equation}\protect\hypertarget{eq:maxent}{}{p(k) = c r^{k},}\label{eq:maxent}\end{equation}

with

\[c = \frac{{1}}{\langle k \rangle-1},\]

and

\[r = \frac{{\langle k \rangle-1}}{\langle k \rangle}.\]

Next, the joint degree distribution \(p(k_{in},k_{out})\) of maximum
entropy was derived using this analytical solution. The joint degree
distribution is the probability of finding a species with \(k_{in}\)
predators and \(k_{out}\) preys in a food web, with
\(k = k_{in} + k_{out}\). We first observed that the joint degree
distribution is identical to the probability of a species having
\(k_{in}\) predators and a total of \(k\) interactions:

\begin{equation}\protect\hypertarget{eq:jointdd}{}{p(k_{in},k_{out}) = p(k_{in},k)}\label{eq:jointdd}\end{equation}

Using Bayes' theorem, we expressed the joint degree distribution as
follows:

\begin{equation}\protect\hypertarget{eq:jointdd2}{}{p(k_{in},k) = p(k_{in}|k)p(k),}\label{eq:jointdd2}\end{equation}

where \(p(k)\) is the degree distribution of maximum entropy derived
above. An unbiased way to estimate \(p(k_{in}|k)\) is by considering all
\(k\) interactions as Bernoulli trials with a probability \(q\) of 0.5
of being incoming interactions. As a result, the joint probability
distribution of maximum entropy of a food web with \(S\) species and a
total of \(L\) interactions is given by:

\begin{equation}\protect\hypertarget{eq:jointdd3}{}{p(k_{in},k) = {k\choose k_{in}}q^{k_{in}}(1-q)^{k-k_{in}}c r^{k},}\label{eq:jointdd3}\end{equation}

which gives, after substitution and simplification:

\begin{equation}\protect\hypertarget{eq:jointdd4}{}{p(k_{in},k_{out}) = {k_{in}+k_{out}\choose k_{in}}c \left(\frac{r}{2}\right)^{k_{in}+k_{out}}}\label{eq:jointdd4}\end{equation}

\hypertarget{predicting-the-adjacency-matrix}{%
\subsubsection{Predicting the adjacency
matrix}\label{predicting-the-adjacency-matrix}}

The number of species \(S\), the number of interactions \(L\), and the
joint degree distribution \(p(k_{in},k_{out})\) are further constraints
on the adjacency matrix of food webs. Directly deriving a network of
maximum entropy considering these constraints is beyond our mathematical
abilities. However, we worked around this issue by first simulating many
networks of \(S\) species whose joint degree distribution was identical
to the one obtained above. We next measured the entropy of all these
potential networks and kept the one with maximum entropy.

Network simulations were conducted in two steps. First, we simulated the
number of species among \(S\) species with \(k_{in}\) in-degrees and
\(k_{out}\) out-degrees using the joint probability distribution (TK).
We removed all occurrences with a total number of interactions different
from \(L\). Next, we generated \emph{a maximum} of 1000 random Boolean
matrices that maintain the numbers of in and out-degrees for all species
in the food web. To do so, we used the curveball algorithm of Strona et
al. (2014) implemented in the Julia library
\texttt{RandomBooleanMatrices} v0.1.1. In order to respect biological
constraints, we removed all occurrences with species having no
interactions

We then computed the entropy of all simulated networks using an entropy
measure for binary matrices (TK). We kept the network with the greatest
entropy value. As a result, for each value of \(S\) we obtained 1000
networks of maximum entropy having different predicted total numbers of
interactions \(L\).

\hypertarget{worldwide-spatial-variation-of-food-web-structure}{%
\subsection{Worldwide spatial variation of food-web
structure}\label{worldwide-spatial-variation-of-food-web-structure}}

A main application of our model is the spatial analysis of food-web
structure at very large spatial scales. We used data from
\href{https://biodiversitymapping.org/}{BiodiversityMapping} to estimate
the number of terrestrial vertebrate species (mammals, birds, and
amphibians) in 10 km x 10 km grids worldwide. For illustrative purposes,
we considered that all species present in a grid were part of the same
food web and interacted with at least one species in their community.

We predicted the structure of all food webs using our model. We first
generated networks of maximum entropy for the range of species richness
encountered in the dataset (TK). We then computed a set of emerging
properties (i.e.~nestedness, maximum trophic level, network diameter,
and entropy) for all of these networks, and mapped the spatial variation
of these measures according to the number of species in each grid. In
figure TK, we present the structure of networks having a number of
interactions \(L\) of maximum a posteriori (MAP) probability.

\hypertarget{data-and-code-availability}{%
\subsection{Data and code
availability}\label{data-and-code-availability}}

All code and data to reproduce this article are available at the Open
Science Framework (TK). Our analyses and simulations were conducted in
Julia v1.5.4.

\hypertarget{results}{%
\section{Results}\label{results}}

Anticipated figures: - Conceptual figure (maybe) - Numerical solutions
as a function of analytical solutions for the maximum entropy degree
distributions - Probability distribution of mean-degree constraints for
a food-web of \(S\) species (maybe) - Maximum entropy distributions of a
food web of \(S\) species and different numbers of links \(L\) - Maxent
against empirical degree distributions - Predicted as a function of
empirical measures of network structure - Predicted and empirical
measures of network structure as a function of species richness - 4
worldwide maps of predicted network properties (nestedness, maximum
trophic level, network diameter, entropy)

\hypertarget{discussion}{%
\section{Discussion}\label{discussion}}

\hypertarget{conclusions}{%
\section{Conclusions}\label{conclusions}}

\hypertarget{acknowledgments}{%
\section{Acknowledgments}\label{acknowledgments}}

This work was supported by the Institute for Data Valorisation (IVADO).

We wrote this article on land located within the traditional unceded
territory of the Saint Lawrence Iroquoian, Anishinabewaki, Mohawk,
Huron-Wendat, and Omàmiwininiwak nations.

\hypertarget{references}{%
\section*{References}\label{references}}
\addcontentsline{toc}{section}{References}

\hypertarget{refs}{}
\begin{CSLReferences}{1}{0}
\leavevmode\hypertarget{ref-Ge2018TurLan}{}%
Ge, Hong, Kai Xu, and Zoubin Ghahramani. 2018. {``Turing: A Language for
Flexible Probabilistic Inference.''} In \emph{International Conference
on Artificial Intelligence and Statistics}, 1682--90. PMLR.
\url{http://proceedings.mlr.press/v84/ge18b.html}.

\leavevmode\hypertarget{ref-MacDonald2020RevLina}{}%
MacDonald, Arthur Andrew Meahan, Francis Banville, and Timothée Poisot.
2020. {``Revisiting the Links-Species Scaling Relationship in Food
Webs.''} \emph{Patterns} 1 (7): 100079.
\url{https://doi.org/10.1016/j.patter.2020.100079}.

\leavevmode\hypertarget{ref-Poisot2016ManMak}{}%
Poisot, Timothee, Benjamin Baiser, Jennifer A. Dunne, Sonia Kefi,
Francois Massol, Nicolas Mouquet, Tamara N. Romanuk, Daniel B. Stouffer,
Spencer A. Wood, and Dominique Gravel. 2016. {``Mangal - Making
Ecological Network Analysis Simple.''} \emph{Ecography} 39 (4): 384--90.
\url{https://doi.org/10.1111/ecog.00976}.

\leavevmode\hypertarget{ref-Strona2014FasUnb}{}%
Strona, Giovanni, Domenico Nappo, Francesco Boccacci, Simone Fattorini,
and Jesus San-Miguel-Ayanz. 2014. {``A Fast and Unbiased Procedure to
Randomize Ecological Binary Matrices with Fixed Row and Column
Totals.''} \emph{Nature Communications} 5 (1, 1): 4114.
\url{https://doi.org/10.1038/ncomms5114}.

\end{CSLReferences}

\end{document}
